\section{System Design}
\label{sec:design}

\subsection{Architecture Overview}
\label{subsec:arch}
Fig.~\ref{fig:arch} depicts the end-to-end deployment. A single AWS
EC2 \texttt{m5.large} instance ($2$\,vCPU, $8$\,GiB RAM) hosts a
lightweight K3s~\cite{rancher2024k3s} v1.30 cluster with three
application namespaces: \texttt{rabbitmq} for the shared broker and
\texttt{baseline}/\texttt{keda} for the two experiment groups. The
\texttt{keda} namespace also hosts the KEDA operator and metrics
adapter installed by Helm. A
Python producer running on the evaluator's workstation injects load
into either \texttt{baseline-queue} or \texttt{keda-queue} over AMQP.
A shell-based collector (\texttt{scripts/collect-metrics.sh}) polls
the RabbitMQ Management API~\cite{broadcom2024rabbitmq} and
\texttt{kubectl} at $\approx$\,$1$~Hz and writes per-second
comma-separated-value (CSV) rows; the same script runs for both
groups so the observation pipeline is identical.

\begin{figure}[t]
  \centering
  % End-to-end architecture. Drawn in TikZ so it scales cleanly in PDF.
\begin{tikzpicture}[
  font=\small,
  box/.style={draw, rounded corners=2pt, align=center, inner sep=3pt,
              minimum height=7mm},
  node distance=3mm and 5mm,
]

% ---------- Laptop (load generator + analysis) ----------
\node[box, fill=blue!8]                         (producer) {Python producer\\{\tiny \texttt{load-test/producer.py}}};
\node[box, fill=blue!8, below=of producer]      (analysis) {Analysis\\{\tiny \texttt{analysis/plot.py}}};
\node[draw, dashed, rounded corners, fit=(producer) (analysis),
      inner sep=4pt, label={[yshift=-2pt]above:\textbf{Laptop}}] (laptop) {};

% ---------- EC2 node ----------
\node[box, fill=orange!15, right=18mm of producer] (rmq)      {RabbitMQ\\{\tiny ns: \texttt{rabbitmq}}};

\node[box, fill=green!10, right=6mm of rmq, yshift=9mm]  (base) {HPA consumer\\{\tiny ns: \texttt{baseline}, cpu=50\%}};
\node[box, fill=red!10,   right=6mm of rmq, yshift=-9mm] (keda) {KEDA consumer\\{\tiny ns: \texttt{keda}, queueLen=100}};

\node[box, fill=gray!10, below=9mm of rmq]     (collect)  {Collector\\{\tiny \texttt{scripts/collect-metrics.sh}}};

\node[draw, dashed, rounded corners,
      fit=(rmq) (base) (keda) (collect),
      inner sep=6pt,
      label={[yshift=-2pt]above:\textbf{EC2 \texttt{m5.large} (K3s v1.30)}}] (ec2) {};

% ---------- Arrows ----------
\draw[-latex] (producer.east) -- ++(0.6,0) |- (rmq.west |- base) node[midway, above,
      font=\scriptsize] {AMQP};
\draw[-latex] (producer.east) -- ++(0.6,0) |- (rmq.west |- keda);

\draw[-latex] (rmq.north east) -- ++(0.4,0) |- (base.west);
\draw[-latex] (rmq.south east) -- ++(0.4,0) |- (keda.west);

\draw[-latex, dashed] (collect.north) -- (rmq.south) node[midway, right,
      font=\scriptsize] {mgmt API};
\draw[-latex, dashed] (collect.east) -| ($(base.south)+(-3mm,0)$);
\draw[-latex, dashed] (collect.east) -| ($(keda.south)+(-3mm,0)$);

\draw[-latex, dashed] (collect.west) -- ++(-5mm,0)
      node[pos=0.6, above, font=\scriptsize] {CSV} |- (analysis.east);

\end{tikzpicture}

  \caption{End-to-end architecture on a single AWS EC2 node. The
  consumer image is shared by both groups; the only variable between
  them is the scaling policy attached to each Deployment.}
  \label{fig:arch}
\end{figure}

\subsection{Consumer Application}
\label{subsec:consumer}
Both groups run the same Spring Boot~\cite{vmware2024springboot}
consumer image (\texttt{htasteqin/cs5296-consumer:v1.0}). Each message
triggers a $\approx$\,$200$\,ms busy-wait so that the workload is
CPU-bound. This is a deliberate design choice: a
\texttt{Thread.sleep} would render HPA blind to the workload---since
the only available signal is CPU utilisation---and would therefore
bias the comparison before it begins. The pod resource settings are
\texttt{requests: 100m~CPU / 256~MiB} and
\texttt{limits: 300m~CPU / 512~MiB}.

\subsection{Scaling Strategies}
\label{subsec:strategies}
The two scalers share identical \texttt{minReplicas}$\,=\,$1 and
\texttt{maxReplicas}$\,=\,$6 (dictated by the 2\,vCPU node), together
with an identical \texttt{behavior} block:

\begin{itemize}[leftmargin=*,itemsep=0pt,topsep=1pt]
  \item \texttt{scaleUp}: \texttt{stabilizationWindowSeconds}\,=\,0;
        policies \texttt{Percent 100\%/15s} $\cup$
        \texttt{Pods 4/15s}, \texttt{selectPolicy}\,=\,\texttt{Max}.
  \item \texttt{scaleDown}: \texttt{stabilizationWindowSeconds}\,=\,60\,s;
        policy \texttt{Percent 50\%/60s}.
\end{itemize}

\noindent Only the scaling \emph{signal} differs:
\textbf{Baseline HPA} uses \texttt{cpu.Utilization} with a
$50\%$ target; \textbf{KEDA} uses a \texttt{rabbitmq}
\texttt{QueueLength} trigger with target~$T\,=\,100$
messages/replica and \texttt{pollingInterval}\,=\,5\,s.

\subsection{Control Variables}
\label{subsec:controls}
To isolate the effect of the scaling signal we hold all non-signal
parameters constant across the two groups: identical Docker image,
identical \texttt{requests}/\texttt{limits}, identical
$\approx$\,$200$\,ms per-message processing cost, a shared RabbitMQ
broker (separate queues per group to prevent interference), the same
EC2 node, and the same burst pattern. The only independent variable
is therefore the scaling signal---CPU utilisation versus queue
length.
